\begin{lemma}[Lewy--Stampacchia order estimate on the full form domain]
Let $F \in L^2(\gamma_n)$ be nonnegative $\gamma_n$-a.e., let $a \in \mathbb{R}$,
and let $u \in \mathcal{K}$ minimize the functional $J$ of
\noderef{ObstacleFunctional} over the cone $\mathcal{K}$ of
\noderef{ObstacleCone}.  Let $\eta$ be the multiplier of
\noderef{ObstacleMultiplier} associated with $u$, $F$, $a$, and put
$k = (a-F)^+ = \max\{a-F, 0\}$.  Then
\[
  \langle \eta, \varphi\rangle \;\leq\; \int k\,\varphi \, d\gamma_n
  \qquad\text{for every } \varphi \in \mathcal{K} .
\]
Thus the multiplier is dominated, in the order of the dual of the form domain, by
the functional induced by the bounded function $(a-F)^+$.  The estimate holds for
every nonnegative form-domain test function, not merely for smooth compactly
supported ones.
\end{lemma}
\begin{proof}
Fix $\varphi \in \mathcal{K}$ and $\varepsilon > 0$.  Put
\[
  z = u - \varepsilon\varphi,
  \qquad z^- = \max\{-z, 0\},
  \qquad v = z^+ = \max\{z,0\} .
\]
Since $u, \varphi \in \mathcal{V}$ and $\mathcal{V}$ is a linear subspace of
$L^2(\gamma_n)$ (its defining condition in \noderef{Energy} is the finiteness of a
weighted $\ell^2$ sum, and the coefficient map of \noderef{HermiteBasis} is
linear, so finiteness is preserved by linear combinations via the triangle
inequality in the weighted $\ell^2$ space), we have $z \in \mathcal{V}$.  The maps
$t \mapsto t^+$ and $t\mapsto t^-$ are $1$-Lipschitz and vanish at $0$, so
\noderef{LatticeContraction} applied to each (with $v = z$ and $w$ the class of
the composite) gives $z^+, z^- \in \mathcal{V}$.  As $z^+ \geq 0$, we get
$v = z^+ \in \mathcal{K}$ by \noderef{ObstacleCone}, so $v$ is an admissible
competitor.

\emph{Step 1: substituting the competitor into the variational inequality.}
Pointwise $z = z^+ - z^-$, and $u = z + \varepsilon\varphi$, so
\[
  v - u = z^+ - (z + \varepsilon\varphi) = z^- - \varepsilon\varphi .
\]
\noderef{VariationalInequality}, applied to the minimizer $u$ and the competitor
$v$, states
$\int (F-a)(v-u)\,d\gamma_n \leq \mathcal{E}(u, v-u)$.  Substituting the display
and using bilinearity of $\mathcal{E}(\cdot,\cdot)$ (\noderef{EnergyBil}) and
linearity of the integral,
\[
  \int (F-a)z^-\,d\gamma_n - \varepsilon\int(F-a)\varphi\, d\gamma_n
    \;\leq\; \mathcal{E}(u, z^-) - \varepsilon\,\mathcal{E}(u,\varphi) .
\]
Rearranging so that the $\varepsilon$-terms stand alone on the left, and using
the definition $\langle\eta,\varphi\rangle = \mathcal{E}(u,\varphi) - \int(F-a)\varphi\,d\gamma_n$
of \noderef{ObstacleMultiplier},
\begin{equation*}
  \varepsilon\,\langle \eta,\varphi\rangle
    \;\leq\; \mathcal{E}(u, z^-) + \int (a-F)\,z^-\, d\gamma_n . \tag{1}
\end{equation*}

\emph{Step 2: the energy term.}
Since $u = z + \varepsilon\varphi$, bilinearity gives
$\mathcal{E}(u, z^-) = \mathcal{E}(z, z^-) + \varepsilon\,\mathcal{E}(\varphi, z^-)$.
Part (2) of \noderef{DisjointCrossTermSign}, applied to $z \in \mathcal{V}$ and
its negative part $z^-$, bounds the first summand by $-\mathcal{E}(z^-)$; and the
Cauchy--Schwarz bound recorded in \noderef{EnergyBil} bounds the second by
$\varepsilon\,\mathcal{E}(\varphi)^{1/2}\mathcal{E}(z^-)^{1/2}$.  Putting
$p = \mathcal{E}(\varphi)^{1/2}$ and $q = \mathcal{E}(z^-)^{1/2}$, both finite
nonnegative reals,
\begin{equation*}
  \mathcal{E}(u,z^-) \;\leq\; -q^2 + \varepsilon p q
    = -\Bigl(q - \frac{\varepsilon p}{2}\Bigr)^2 + \frac{\varepsilon^2 p^2}{4}
    \;\leq\; \frac{\varepsilon^2}{4}\,\mathcal{E}(\varphi) . \tag{2}
\end{equation*}

\emph{Step 3: the load term.}
Because $u \geq 0$ and $\varphi \geq 0$ $\gamma_n$-a.e.\ (both lie in
$\mathcal{K}$, see \noderef{ObstacleCone}), we have pointwise a.e.
\[
  z^- = \max\{\varepsilon\varphi - u,\, 0\} \leq \varepsilon\varphi,
  \qquad z^- \geq 0 .
\]
Now $k = (a-F)^+$ satisfies $a - F \leq k$ and $k \geq 0$ pointwise, so
multiplying the first inequality by the nonnegative number $z^-$ and then using
$z^- \leq \varepsilon\varphi$ and $k \geq 0$,
\[
  (a-F)\,z^- \;\leq\; k\,z^- \;\leq\; \varepsilon\, k\,\varphi
  \qquad\gamma_n\text{-a.e.}
\]
Integrating,
\begin{equation*}
  \int (a-F)\,z^-\, d\gamma_n \;\leq\; \varepsilon \int k\,\varphi\, d\gamma_n .
  \tag{3}
\end{equation*}
The right-hand integral is finite: since $F \geq 0$ a.e.\ we have
$0 \leq k = (a-F)^+ \leq \max\{a,0\}$, so $k$ is bounded, and
$\varphi \in L^2(\gamma_n) \subseteq L^1(\gamma_n)$ by \noderef{L2IntoL1}.  The
left-hand integral is finite for the same reason, $z^-$ being in
$L^2(\gamma_n)$.

\emph{Step 4: conclusion.}
Combining (1), (2) and (3),
\[
  \varepsilon\,\langle\eta,\varphi\rangle
    \;\leq\; \frac{\varepsilon^2}{4}\mathcal{E}(\varphi)
             + \varepsilon\int k\varphi\, d\gamma_n .
\]
Dividing by $\varepsilon > 0$,
\[
  \langle\eta,\varphi\rangle
    \;\leq\; \frac{\varepsilon}{4}\mathcal{E}(\varphi) + \int k\varphi\, d\gamma_n .
\]
The test function $\varphi$ is fixed with $\mathcal{E}(\varphi) < \infty$, and
$\varepsilon > 0$ was arbitrary, so letting $\varepsilon \downarrow 0$ removes the
first term and gives
$\langle\eta,\varphi\rangle \leq \int k\varphi\, d\gamma_n$.
\end{proof}
